% from ega0-22-fr.tex
\oldpage[0\textsubscript{IV-1}]{199}

$(a_i)_{s+1\leqslant i\leqslant r}$ est $p$-libre sur $K^p$ et forme une
$p$-base sur $K^p$ du corps $K^p(a_1,\ldots,a_r)$ de sorte que l'on a
\[
  a_i\in K^p(a_{s+1},\ldots,a_r)\qquad\text{pour }1\leqslant i\leqslant s.
\]
Désignons par $A''$ l'anneau quotient de $A[T_{s+1},\ldots,T_r]$ par l'idéal
engendré par les $F_i$ d'indice $i\geqslant s+1$~; alors, par (22.4.4), $A''$
est un anneau local dont l'idéal maximal est $\mathfrak m''=\mathfrak mA''$
(donc de carré nul) et le corps résiduel
$K''=K(a_{s+1}^{1/p},\ldots,a_r^{1/p})$. On a donc
$a_i\in K''^p$ pour $1\leqslant i\leqslant s$ et par suite, dans $A''$,
l'élément $f_i$, pour $1\leqslant i\leqslant s$, peut s'écrire
\begin{equation}
\label{0.22.4.7.3-fr}
  f_i=g_i^p+h_i
  \tag{22.4.7.3}
\end{equation}
où $g_i\in A''$ et $h_i\in\mathfrak m''$ (comme $\mathfrak m''^2=0$ et
$p\mathfrak m''=0$, il est immédiat que $h_i$ est déterminé de façon
\emph{unique} par ces conditions). Remplaçant $T_i$ par $T_i+g_i$, on voit donc
que si l'on pose
\[
  G_i(T_i)=T_i^p+\sum_{k=1}^{p-1}\binom pk g_i^{p-k}T_i^k-h_i
  \qquad(1\leqslant i\leqslant s),
\]
$B$ est l'anneau quotient de $A''[T_1,\ldots,T_s]$ par l'idéal engendré par les
$G_i$~; or, tous les coefficients de $G_i$ sauf le coefficient dominant
appartiennent à $\mathfrak m''$, donc on est dans la situation de (22.4.3),
$B$ est un anneau local, et son corps résiduel $K'$ est \emph{isomorphe} à
$K''$, ce qui prouve (i) (qui d'ailleurs résulte directement de ce que
$B^p\subset A$, et par suite il n'y a qu'un seul idéal de $B$ au-dessus de
$\mathfrak m$).

Notons maintenant que pour que (22.4.7.2) soit injectif, il faut et il suffit
qu'il en soit ainsi de
\[
  \pi'_{B_{(K)}/K}:\Omega^1_{B_{(K)}/K}\otimes_{B_{(K)}}K
  \longrightarrow\Omega^1_A\otimes_AK
\]
(par (22.4.6)). Or $B_{(K)}=B/\mathfrak mB$ est l'algèbre quotient de
$C=K[T_1,\ldots,T_r]$ par l'idéal engendré par les polynômes
$\overline F_i=T_i^p-a_i$, et comme $d_{C/K}\overline F_i=0$ il résulte de
(20.5.13) que $\Omega^1_{B_{(K)}/K}$ est un $B_{(K)}$-module libre ayant pour
\emph{base} les $d_{B_{(K)}/K}t_i$, où $t_i$ est l'image canonique de $T_i$
$(1\leqslant i\leqslant r)$. Mais par définition (22.4.5) l'image canonique par
$j':B_{(K)}\longrightarrow A$ d'une classe mod.~$\mathfrak mB$ d'un élément
$x\in B$ est l'élément $x^p\in A$~; on en déduit aussitôt, puisque $t_i^p=f_i$,
que l'image de $d_{B_{(K)}/K}t_i\otimes1$ par $\pi'_{B_{(K)}/K}$ est
$d_Af_i\otimes1$ dans $\Omega^1_A\otimes_AK$~; la condition que
$\pi'_{B_{(K)}/K}$ soit injectif est donc équivalente au fait que les
$d_Af_i\otimes1_K$ sont \emph{linéairement indépendants sur $K$}, ou encore
que leurs images $d_Af_i\otimes1_{K'}$ sont \emph{linéairement indépendants
sur $K'$} dans $\Omega^1_A\otimes_AK'$.

Or, si l'on applique (22.4.4) à la $A$-algèbre $A''$, on voit que le noyau de
l'homomorphisme canonique
$\Omega^1_A\otimes_AK'\longrightarrow\Omega^1_{A''}\otimes_{A''}K'$ a pour
base les $d_Af_i\otimes1_{K'}$ pour $s+1\leqslant i\leqslant r$. La condition
précédente équivaut donc encore à dire que les $d_{A''}f_i\otimes1_{K'}$ sont
linéairement indépendants dans $\Omega^1_{A''}\otimes_{A''}K'$ pour
$1\leqslant i\leqslant s$. Or, on a $d_{A''}f_i=d_{A''}h_i$ par (22.4.7.3)~;
d'autre part, si $V''$ est le $K'$-espace vectoriel $\mathfrak m''$, on a
$\operatorname{rg}_{K'}(V'')=\operatorname{rg}_K(V)$ (22.4.1, (ii)), et la condition
$\operatorname{rg}_{K'}(V')=\operatorname{rg}_K(V)$ est donc équivalente à
$\operatorname{rg}_{K'}(V'')=\operatorname{rg}_{K'}(V')$. Mais il résulte de (22.4.3) que cette dernière
relation est équivalente au fait que les classes $h_i$

